%% CPP 2027 Submission — 0pat Exercise XXIX: Observing the Intersections of the Complex Plane's Four Circles
\documentclass[sigplan,10pt,anonymous,review]{acmart}
\settopmatter{printfolios=true,printacmref=false}

\AtBeginDocument{%
  \providecommand\BibTeX{{\normalfont B\scshape ib\TeX}}}

\settopmatter{printacmref=false}
\renewcommand\footnotetextcopyrightpermission[1]{}
\pagestyle{plain}

\usepackage{microtype}
\usepackage{booktabs}
\usepackage{amsmath}
% XeTeX (tectonic): acmart loads newtxmath/libertine which already defines \Bbbk;
% reset it so amssymb's definition is accepted.
\let\Bbbk\relax
\usepackage{amssymb}

\begin{document}

\title{Observing the Intersections of the Complex Plane's Four Objects\\
(Exercise XXIX)}

\author{Anonymous}
\affiliation{%
  \institution{Anonymous}
}
\email{anonymous@example.org}

\begin{abstract}
We observe, directly and with machine-checked proofs, the intersections
of the four central objects of the Riemann direction: the real axis, the
imaginary axis, the critical circle, and the prime circles. The complete
intersection table is established: the real axis meets the critical
circle at $\{0, 2\}$ and the imaginary axis at $\{0\}$; the real axis
meets the prime circle of $p$ at $\{\pm\sqrt{p}\}$ and the imaginary
axis at $\{\pm i\sqrt{p}\}$; the prime-2 circle meets the critical
circle at $\{1 \pm i\}$; the prime-3 circle meets it at
$\{3/2 \pm \sqrt{3}/2 \cdot i\}$; and --- the new structural theorem ---
the prime circle of $p \ge 5$ does not meet the critical circle at all:
a point of both circles would force $a = p/2$ and
$b^2 = p(4-p)/4 < 0$. Equivalently: the prime circle (radius
$\sqrt{p}$, center $0$) meets the critical circle (center $(1,0)$,
radius $1$) if and only if $p \in \{2, 3\}$. Five theorems, zero
errors, zero \texttt{sorry}.
\end{abstract}

\keywords{Lean, formalization, critical circle, prime circle, intersections, 0pat}

\maketitle

\section{The observation}

The Riemann direction has four central geometric objects in the complex
plane: the real axis, the imaginary axis, the critical circle (center
$(1,0)$, radius $1$, the image of the critical line under reciprocal),
and the prime circle of each prime $p$ (center $0$, radius $\sqrt{p}$,
the locus of the two-square representations of $p$). This exercise
observes their intersections --- the complete table, every entry
machine-checked.

\begin{center}
\begin{tabular}{lll}
\toprule
Pair & Intersection & Status \\
\midrule
real axis $\cap$ imaginary axis & $\{0\}$ & definition \\
real axis $\cap$ critical circle & $\{0, 2\}$ & C020 (proved) \\
imaginary axis $\cap$ critical circle & $\{0\}$ & new \\
real axis $\cap$ prime circle $p$ & $\{\pm\sqrt{p}\}$ & new \\
imaginary axis $\cap$ prime circle $p$ & $\{\pm i\sqrt{p}\}$ & new \\
prime-2 circle $\cap$ critical circle & $\{1 \pm i\}$ & C020 (proved) \\
prime-3 circle $\cap$ critical circle & $\{3/2 \pm \sqrt{3}/2 \cdot i\}$ & new \\
prime circle $p \ge 5$ $\cap$ critical circle & $\emptyset$ & new (structural) \\
\bottomrule
\end{tabular}
\end{center}

\section{The theorems}

\paragraph{The imaginary axis meets the critical circle only at the
origin (imagAxis\_inter\_criticalCircle).} A point $t \cdot i$ of the
imaginary axis has distance $\sqrt{1 + t^2}$ from $(1,0)$; it lies on
the critical circle exactly when $1 + t^2 = 1$, i.e. $t = 0$. The
contrast with the real axis is exact: the real axis meets the critical
circle at the two points $\{0, 2\}$ (C020), the imaginary axis at one
point $\{0\}$.

\paragraph{The axes meet the prime circle in square roots
(realAxis\_inter\_primeCircle, imagAxis\_inter\_primeCircle).}
A real point $r$ has norm $r^2$; it lies on the prime circle of $p$
exactly when $r^2 = p$, i.e. $r = \pm\sqrt{p}$. Similarly, the
imaginary point $t \cdot i$ lies on the prime circle exactly when
$t^2 = p$. Both intersections are irrational pairs for prime $p$;
they are never lattice points.

\paragraph{The prime-3 circle meets the critical circle
(prime3Circle\_inter\_criticalCircle).} The point
$\frac32 + \frac{\sqrt{3}}{2} i$ has norm $9/4 + 3/4 = 3$ and distance
$\sqrt{1/4 + 3/4} = 1$ from $(1,0)$. Together with the prime-2
intersection $\{1 \pm i\}$ (C020): the small primes $2$ and $3$ have
circles that meet the critical circle.

\paragraph{The structural theorem: circles of primes $p \ge 5$ do not
meet the critical circle (primeCircle\_inter\_criticalCircle\_ge5).}
Suppose $z = a + bi$ lies on both circles. Then $\|z\|^2 = p$ and
$\|z - 1\|^2 = 1$. Subtracting gives $-2a + 1 = 1 - p$, so $a = p/2$;
then $b^2 = p - a^2 = p - p^2/4 = p(4-p)/4$, which is negative for
$p \ge 5$ --- contradicting $b^2 \ge 0$. Hence the prime circle of
$p \ge 5$ is disjoint from the critical circle. Equivalently: the prime
circle meets the critical circle if and only if $p \in \{2, 3\}$.

\section{Honest boundary and position}

The content is classical geometry (circle intersections in the plane);
the value is observational and structural. The observation completes the
intersection table of the Riemann direction's four objects and adds one
new structural fact: only the prime circles of $2$ and $3$ touch the
critical circle; all other prime circles are disjoint from it. No claim
is made about the zeros of $\zeta$: the critical circle is the geometric
image of the critical line, and its intersections with the prime circles
describe the geometry of the plane, not the analysis of the function.

\section{Formalization}

The proof is a single Lean file (\texttt{proof.lean}, 108 lines),
importing \texttt{Formal.ZeroRelative.ComplexAxis} (the base
construction of the Riemann chain, C011) and
\texttt{Mathlib.Analysis.Real.Sqrt}. Five theorems, compiled with Lean
4.32.2 / mathlib v4.32.2, zero errors, zero \texttt{sorry}.

\begin{center}
\begin{tabular}{lll}
\toprule
Theorem & Statement & Proof core \\
\midrule
\texttt{imagAxis\_inter\_criticalCircle} & $t \cdot i \in$ crit. circle $\leftrightarrow t = 0$ & simp; nlinarith \\
\texttt{realAxis\_inter\_primeCircle} & $r \in$ circle$_p \leftrightarrow r^2 = p$ & simp \\
\texttt{imagAxis\_inter\_primeCircle} & $t \cdot i \in$ circle$_p \leftrightarrow t^2 = p$ & simp \\
\texttt{prime3Circle\_inter\_criticalCircle} & $\frac32 + \frac{\sqrt3}{2}i$ in both & nlinarith \\
\texttt{primeCircle\_inter\_criticalCircle\_ge5} & $p \ge 5 \Rightarrow$ disjoint & subtraction; nlinarith \\
\bottomrule
\end{tabular}
\end{center}

\section{Conclusion}

The four objects of the Riemann direction are now fully observed: their
intersection table is complete and machine-checked, and the new
structural fact stands --- the prime circle touches the critical circle
if and only if the prime is $2$ or $3$; for $p \ge 5$ the circles are
disjoint. The geometry of the plane is now mapped exactly; the analysis
that would connect it to the zeros of $\zeta$ remains the open span of
the bridge.

\bibliographystyle{ACM-Reference-Format}
\bibliography{refs}

\end{document}
